Данная работа посвящена исследованию ресурсных систем массового обслуживания с потерями, в которых каждая поступающая заявка занимает один обслуживающий прибор и случайный объём ограниченного ресурса (модель $GI/G/N/N/R$). В отличие от классических марковских моделей, в фокусе исследования находится анализ чувствительности стационарных характеристик системы к форме распределения входящего потока и процесса обслуживания при немарковских предположениях.

Ввиду вычислительной сложности точного аналитического расчета подобных многомерных систем, в работе разработан высокопроизводительный программный комплекс для дискретно-событийного имитационного моделирования. Выявлено влияние вариации межприходных интервалов на вероятность блокировки заявок и исследованы границы применимости свойства инвариантности при сохранении фиксированных первых моментов случайных величин.
